\chapter{Verification and Validation}
\label{ch:verification}

The single-leg CFD provides all quantitative statements of \cref{ch:results}. This chapter checks it in three
ways: against a flapping-foil experiment (validation), on a coarser mesh (spatial discretization) and from cycle
to cycle (temporal convergence). It closes with the consequences for the drag-type cases, which were not part of
the mesh study, and with a check of the real-time solver Flare Live against the CFD.

\section{Validation with a Heaving and Pitching Foil}
\label{sec:ver-v1}

\paragraph{Case.} Schouveiler et al.\ \cite{Schouveiler2005} measured the thrust and efficiency of a NACA\,0012
foil in combined heave and pitch in a towing tank, with error bars. One of their points is reproduced with the
same solver, the same overset method and comparable resolution as the fluke cases: chord $c=\SI{0.1}{\metre}$,
span \SI{0.6}{\metre} with free tips, heave amplitude $h_0/c=0.75$, pitch axis at $c/3$, pitch leading heave by
\SI{90}{\degree}, $\mathit{St}=0.25$, maximum angle of attack \SI{15}{\degree}, $\U=\SI{0.4}{\metre\per\second}$
($f=\SI{0.667}{\hertz}$) and $\mathit{Re}=4\times10^4$. The component mesh has \SI{2.5}{\milli\metre} cells
($c/40$), the background \SI{5}{\milli\metre}. Three cycles were simulated without a start-up ramp.

\paragraph{Results.} The streamwise force settles after the first cycle (\cref{fig:verification}a); cycles 2 and 3
differ by less than 0.1\,\%. The simulated motion reaches $\alpha_{75}=\SI{14.6}{\degree}$, consistent with the
prescribed \SI{15}{\degree}, and the mean lateral force is 1.3\,\% of the thrust. The mean thrust is
\SI{1415}{\milli\newton} and the input power \SI{946}{\milli\watt}. \Cref{tab:v1} compares the coefficients
with the experiment.

\begin{table}[t]
  \centering
  \caption[Validation against Schouveiler et al.]{Validation against the experiment of Schouveiler et
    al.\ \cite{Schouveiler2005} (NACA\,0012, $\mathit{St}=0.25$, $\alpha_\text{max}=\SI{15}{\degree}$).}
  \label{tab:v1}
  \small
  \begin{tabular}{@{}lccc@{}}
    \toprule
    & CFD (this work) & Experiment & Deviation \\
    \midrule
    Thrust coefficient $C_T$ & 0.295 & $0.32\pm0.01$ & $-8$\,\% \\
    Froude efficiency $\etaF$ & 0.60 & $0.73\pm0.04$ & $-18$\,\% \\
    \bottomrule
  \end{tabular}
\end{table}

The thrust coefficient is within 8\,\% of the measurement, inside the $\pm$10--15\,\% usually reported for
flapping-foil CFD, and is taken as the main criterion. The efficiency is 18\,\% too low. Three causes act in the
direction observed. (i) The simulated foil has free tips and an aspect ratio of 6, while the experiment was close
to two-dimensional; a finite-wing estimate $\mathit{AR}/(\mathit{AR}+2)=0.75$ suggests tip losses of 10--20\,\%, which
affect efficiency more than thrust because tip vortices cost power. (ii) At $\mathit{Re}=4\times10^4$ the boundary
layer is transitional; a laminar model separates earlier and overestimates the power. (iii) The mesh ($c/40$) is
coarser than in the fluke cases ($c/50$). Because this thesis compares propulsors simulated with the same solver
and settings, such systematic deviations affect both propulsors in the same direction to first order; the tip losses and the laminar separation act mainly on the lift-based fluke, so its advantage is, if anything, underestimated. Absolute efficiencies should be read as conservative.

\begin{figure}[t]
  \centering
  \includegraphics[width=\textwidth]{fig_verification}
  \caption[Verification and validation]{Verification and validation of the single-leg CFD. (a) Validation case:
    streamwise force coefficient with cycle means (dashed). (b) Thrust coefficient and efficiency against the
    experiment of Schouveiler et al.\ \cite{Schouveiler2005}. (c) Fluke L0 on the fine and the coarse mesh over the
    common time window. (d) Change of the mean thrust between the last and the second-last cycle for all fluke cases
    and the validation case.}
  \label{fig:verification}
\end{figure}

\section{Mesh Sensitivity}
\label{sec:ver-mesh}

The reference fluke case L0 was repeated on a mesh in which all cell sizes were multiplied by $r=1.33$ (component
\SI{1.6}{\milli\metre}, background \SI{3.3}{\milli\metre}). The coarse run was stopped at $t=\SI{0.91}{\second}$,
so both runs are compared over the common window $t=\SIrange{0.41}{0.91}{\second}$. The force histories are
nearly identical in shape (\cref{fig:verification}c). The coarse mesh gives 7.3\,\% less thrust
(\SI{19.75}{} versus \SI{21.30}{\milli\newton}), 3.9\,\% less power and 3.4\,\% lower efficiency.

With two meshes the order of convergence cannot be computed. Assuming second order ($p=2$) and a safety factor of
$F_s=1.25$ \cite{Roache1994}, the grid convergence index of the thrust on the fine mesh is
\begin{equation}
  \mathrm{GCI} = F_s\,\frac{|\varepsilon|}{r^p-1} = 1.25\times\frac{0.073}{1.33^2-1} \approx 12\,\%.
  \label{eq:gci}
\end{equation}
The discretization uncertainty of the fluke thrust is therefore taken as about $\pm12$\,\%, and that of power and
efficiency as about $\pm6$\,\%. Since the coarse mesh gives lower thrust, the fine-mesh values are more likely to
be slightly low than high.

\section{Temporal Convergence}
\label{sec:ver-time}

For every fluke case, the mean thrust of the last cycle differs from that of the second-last cycle by at most
0.61\,\% (L2 at $\U=0$), and by at most 0.31\,\% for all other cases (\cref{fig:verification}d). The validation
case differs by less than 0.1\,\%. Averaging over the last two cycles after a half-cycle start-up is therefore
sufficient; the periodic state is reached within the first full cycle. The fan-paddle cases were evaluated in their
third cycle. For the robot gait at \SI{0.08}{\metre\per\second}, cycles 2 and 3 give \SI{7.3}{} and
\SI{5.2}{\milli\newton}; this case, whose thrust is small, has a correspondingly larger uncertainty.

\section{Consequences for the Drag-Type Cases}
\label{sec:ver-drag}

No separate mesh study was run for the fan paddle, and its mesh is coarser than that of the fluke
(\cref{tab:meshes}). Three arguments bound the effect on the conclusions.
\begin{enumerate}
  \item \emph{Flow physics.} The paddle is a thin plate with sharp edges moved broadside. Its separation lines are
        fixed at the edges, whereas the fluke's thrust depends on a leading-edge vortex and leading-edge suction
        that are more sensitive to resolution. The fluke's $\pm12$\,\% is a plausible upper bound for the paddle.
  \item \emph{Low speed.} Up to \SI{0.15}{\metre\per\second} the optimized paddle produces 1.6--2.0 times the thrust
        of the fluke. An error of $\pm12$\,\% does not change this ranking.
  \item \emph{Crossover.} Near the crossover the thrust falls with slopes of about \SI{-110}{} (fluke) and
        \SI{-340}{\milli\newton\second\per\metre} (paddle). A $\pm12$\,\% error of the paddle thrust
        (about \SI{2.5}{\milli\newton}) shifts the crossover by about \SI{0.01}{\metre\per\second}; opposite
        errors of $\pm12$\,\% on both sides shift it by about \SI{0.02}{\metre\per\second}. The uncertainty of the
        added-mass correction (\cref{sec:cfd-composite}) adds less than \SI{0.01}{\metre\per\second}.
\end{enumerate}
The coarser paddle mesh is nevertheless a limitation and is listed in \cref{sec:disc-limits}.

\section{Cross-Check of the Real-Time Solver}
\label{sec:ver-flare}
\label{sec:res-flare}

Flare Live (\cref{sec:cfd-flare}) simulates the whole robot in real time but on a lattice of \SIrange{2}{3}{\milli\metre}. Before
any of its results can be used, its fluke force is checked against the CFD. In the tow runs, the fluke performs exactly the L0 motion of the CFD. Its
angle of attack agrees ($\alpha_{75}=\SI{19.7}{\degree}$ versus \SI{19.9}{\degree} at \SI{0.223}{\metre\per\second}), but its
thrust does not (\cref{fig:flare-fin}). On the \SI{3}{\milli\metre} lattice, Flare Live gives \SI{20.0}{},
\SI{8.5}{} and \SI{-0.1}{\milli\newton} at \SIlist{0;0.1;0.223}{\metre\per\second}, against \SI{45.5}{}, about \SI{34}{} (interpolated) and \SI{21.4}{\milli\newton} in CFD (\cref{tab:app-flare}): 44\,\%, 25\,\% and practically zero. The deficit is about
\SIrange{21}{26}{\milli\newton} at all three speeds, which is the order of the lift-based thrust itself.

The idle fluke, held at its zero position, has a drag of \SI{12.6}{\milli\newton} at \SI{0.223}{\metre\per\second}. At its
incidence of \SI{-25}{\degree} this corresponds to a drag coefficient of about 0.45 based on the planform, a normal value for a
stalled plate. The thrust increment of flapping over idling is \SI{12.5}{\milli\newton}, still well below the CFD thrust.

On a \SI{2}{\milli\metre} lattice the fluke thrust at \SI{0.223}{\metre\per\second} rises to \SI{4.1}{\milli\newton} (19\,\% of
CFD) and the idle drag hardly changes (\SI{11.6}{\milli\newton}). The result is not converged, and a Richardson extrapolation from
the two lattices gives only about \SIrange{8}{12}{\milli\newton}, which indicates that the lattices are far from the asymptotic
range. The cause is resolution: the \SI{34}{\milli\metre} chord spans only 11 cells at \SI{3}{\milli\metre} and 17 at
\SI{2}{\milli\metre}, while the boundary layer at $\mathit{Re}\approx7600$ is about \SI{0.4}{\milli\metre} thick. The leading edge,
where the suction that drives lift-based thrust acts, is not resolved. Resolving it would need a lattice of \SI{1}{\milli\metre}
or finer at many times the cost, which defeats the purpose of a real-time solver.

\begin{figure}[t]
  \centering
  \includegraphics[width=\textwidth]{fig_flare_fin}
  \caption[Flare Live compared with CFD]{Fluke force in Flare Live and in CFD for the same fluke and motion. (a) Streamwise force
    versus tow speed, flapping and idle. (b) Ratio of the Flare to the CFD thrust.}
  \label{fig:flare-fin}
\end{figure}

Flare Live is therefore used only for the reference motion of the fluke and for qualitative trends of the free-swimming
robot, which are collected in \cref{app:flare-trends}; all quantitative statements rest on the CFD.

\section{Summary}
\label{sec:ver-summary}

\Cref{tab:uncertainty} collects the checks of this chapter.

\begin{table}[H]
  \centering
  \caption[Uncertainty summary]{Summary of the verification and validation.}
  \label{tab:uncertainty}
  \small
  \begin{tabularx}{\textwidth}{@{}l>{\raggedright\arraybackslash}X@{}}
    \toprule
    Check & Result \\
    \midrule
    Validation (foil experiment) & $C_T$ $-8$\,\%, $\etaF$ $-18$\,\%; deviations mainly from free tips and laminar model \\
    Mesh sensitivity (fluke) & $\Tm$ $-7.3$\,\% on the $\times1.33$ coarser mesh, GCI $\approx\pm12$\,\% \\
    Temporal convergence & fluke: last two cycles differ by $\le0.61$\,\%; robot gait at \SI{0.08}{\metre\per\second}: 7.3 vs.\ \SI{5.2}{\milli\newton} \\
    Drag-type mesh & not studied; ranking insensitive, crossover within $\pm\SI{0.02}{\metre\per\second}$ \\
    Real-time solver (Flare Live) & fluke thrust 0--44\,\% of CFD on 2--3\,mm lattices; used for trends only \\
    \bottomrule
  \end{tabularx}
\end{table}

The single-leg CFD reproduces the thrust of a flapping foil well and underestimates its efficiency. Its numerical
uncertainty is dominated by spatial discretization, about $\pm12$\,\% in thrust. Differences between propulsors of
less than about 15\,\% are therefore not interpreted as significant in \cref{ch:results}.
